\bookchapter{Capítulo 39: Engenharia de contexto — fazer a IA entender a sua intenção}

\chapterintro{Já aconteceu de trocar de conversa e a IA “esquecer” tudo do seu projeto? Não é de propósito: ela é naturalmente esquecida. Este capítulo ensina a instalar um “sistema de memória de longo prazo” na IA e apresenta um dos paradigmas mais quentes de 2026 — o desenvolvimento orientado a especificação.}

\booksection{Arquivo de memória do projeto: a “memória de longo prazo” da IA}

Os modelos de IA têm uma limitação fundamental: a memória é de curto prazo. A cada nova conversa, a IA começa do zero sobre o seu projeto. Você precisa repetir “usamos React”, “o backend é Node.js”, “o banco é PostgreSQL”, “o estilo usa aspas duplas”. É ineficiente e fácil de esquecer — como treinar um estagiário novo todo dia.

A solução é o “arquivo de memória do projeto”: um texto na raiz do projeto que a ferramenta de IA lê automaticamente ao iniciar a conversa. Cada ferramenta tem o seu nome:

\begin{center}
\begin{tabularx}{\linewidth}{|X|X|}
\hline
\textbf{Ferramenta} & \textbf{Nome do arquivo} \tabularnewline
\hline
Cursor & \texttt{.\allowbreak cursorrules} \tabularnewline
Claude Code & \texttt{CLAUDE.\allowbreak md} \tabularnewline
GitHub Copilot & \texttt{copilot-\allowbreak instructions.\allowbreak md} \tabularnewline
Solução genérica & \texttt{AGENTS.\allowbreak md} \tabularnewline
\hline
\end{tabularx}
\end{center}

\begin{center}\small\color{gray}Convenções de nome do arquivo de memória por ferramenta\end{center}

Os nomes mudam, o princípio é o mesmo: escreva no arquivo o que a IA deve “sempre lembrar”. Um bom arquivo traz: visão geral do projeto (o que é), stack (tecnologias), padrões de código (nomenclatura, indentação, comentários), estrutura de pastas e convenções (branches, commits, testes). O CLAUDE.md de Karpathy tem só 70 linhas, mas cobre o essencial — conciso e completo.

\begin{lstlisting}
# Projeto: app pessoal de controle financeiro
## Visão geral
Um app full-stack de finanças pessoais que ajuda a registrar receitas e despesas do dia a dia.
## Stack tecnológica
- Frontend: React + TypeScript + Tailwind CSS
- Backend: Node.js + Express + TypeScript
- Banco de dados: SQLite (dev) / PostgreSQL (produção)
- Testes: Vitest
## Padrões de código
- Usar aspas duplas
- Indentação de 2 espaços
- Nomes de função em camelCase
- Definições de tipo em PascalCase
- Toda função deve ter comentário JSDoc
## Estrutura de diretórios
/src
  /frontend  - código frontend
  /backend   - código backend
  /shared    - tipos compartilhados entre frontend e backend
  /tests     - arquivos de teste
## Convenções de desenvolvimento
- Mensagens de commit em português
- Recursos novos exigem testes primeiro
- APIs sempre retornam { success, data, error }
\end{lstlisting}

\begin{tcbwarn}{Atenção: não vire um "segundo documento de requisitos"}
Cuidado com o tamanho: não escreva dezenas de páginas. O objetivo é “dar contexto rápido à IA”, não “descrever cada função em detalhe”. Mantenha curto, estruturado e só com “o que a IA precisa saber”.

\end{tcbwarn}

\booksection{Desenvolvimento orientado a especificação (SDD)}

O Spec-Driven Development (SDD) amadureceu em 2026, com uma ideia em uma frase: escreva a especificação antes do código. Diferente do TDD clássico, a “spec” do SDD não é só casos de teste: inclui requisitos, definição de interfaces, modelos de dados etc. A IA gera o código a partir dela; os humanos escrevem e mantêm a especificação.

Fluxo do SDD em cinco passos:

\begin{enumerate}
  \item Descreva os requisitos em linguagem natural (em um arquivo .md);
  \item Defina interfaces e estruturas de dados (no mesmo arquivo);
  \item Peça à IA que gere a implementação a partir da spec;
  \item Rode testes para validar se a entrega segue a spec;
  \item Se não seguir, corrija a spec ou peça à IA para corrigir o código.
\end{enumerate}

Nesse fluxo, o arquivo de especificação é a “fonte única da verdade”: o código pode ser regenerado, mas a spec precisa estar correta primeiro.

Escrever o arquivo de especificações

A especificação está completa?

Completa

A IA gera código a partir da especificação

Rodar testes para verificar o resultado

Incompleta / testes falham

Problema no código corrige o código; problema na especificação corrige a especificação

De volta à verificação

Tudo aprovado entregar

\begin{center}\small\color{gray}Figura 3: fluxo de trabalho do desenvolvimento orientado a especificação (SDD)\end{center}

A vantagem do SDD é que a spec é mais estável que o código. Quando o modelo ou a ferramenta muda, a sua spec continua valendo — basta gerar o código de novo. É por isso que o livro insiste que “metodologia vale mais que ferramenta”: a spec é o veículo concreto da sua metodologia.

\booksection{Gestão da janela de contexto}

Mesmo com o arquivo de memória, ainda faltará contexto em conversas longas. A janela de contexto (quanto a IA processa ao mesmo tempo) é limitada. Com chat muito extenso e muitos arquivos citados, a IA “esquece” o começo e a qualidade cai — como uma conversa de seis horas em que o amigo já não lembra o que você disse na abertura.

A estratégia central é a “referência em camadas”:

\begin{center}
\begin{tabularx}{\linewidth}{|X|X|X|}
\hline
\textbf{Camada de contexto} & \textbf{Conteúdo} & \textbf{Como gerenciar} \tabularnewline
\hline
1ª · Memória do projeto & Stack, convenções, estrutura de pastas & Em CLAUDE.md etc., carregado automaticamente, sem repetir \tabularnewline
2ª · Referência de arquivos & Código envolvido na tarefa atual & Cite com @ com precisão; não cole o arquivo inteiro no chat \tabularnewline
3ª · Contexto da conversa & Histórico do chat da tarefa atual & Ao concluir uma tarefa, abra um chat novo e evite inflar o contexto \tabularnewline
4ª · Instrução imediata & O que se pede neste passo & Clara, curta e objetiva; sem descrições intermináveis \tabularnewline
\hline
\end{tabularx}
\end{center}

\begin{center}\small\color{gray}Estratégia de gestão da janela de contexto em camadas\end{center}

\booksection{Pontos práticos de engenharia de contexto}

Engenharia de contexto não é uma configuração única; é um processo de melhoria contínua. Alguns pontos comprovados na prática:

\begin{enumerate}
  \item Atualize o arquivo de memória junto com o projeto. Ao adotar uma tecnologia nova ou mudar a arquitetura, atualize o arquivo na hora — não explique de improviso no chat.
  \item Prefira @ a colar textões. A maioria das ferramentas aceita @arquivo ou @pasta, injetando o conteúdo no contexto em segundo plano — mais eficiente e sem poluir a conversa.
  \item Conversas curtas e cirúrgicas. Cada chat, uma tarefa ou um ponto de função. Ao concluir, abra outro. Conversas longas gastam a janela de contexto e "distraem" a IA.
  \item Mostre o resultado em vez de descrever. Em bugs, @cite o log de erro em vez de dizer "parece que tem algum erro". A IA responde melhor à informação original do que à sua versão.
  \item Limpe e arquive com regularidade. Conversas concluídas e arquivos temporários ocupam o espaço de trabalho. Manter o diretório limpo também ajuda a IA a entender a estrutura.
\end{enumerate}

\begin{tcbpractice}{Mão na massa: escreva um arquivo de memória de 70 linhas}
Escolha um projeto em andamento (ou pretendido) e escreva um CLAUDE.md: visão geral (2–3 frases), stack (liste), padrões de código (3–5), estrutura de pastas (árvore) e convenções (2–3). Fique em até 70 linhas — ao final, a sua visão do projeto já estará mais clara.

\end{tcbpractice}

\begin{tcbsummary}{Resumo do capítulo}
\begin{itemize}
  \item A memória da IA é de curto prazo: a cada conversa, começa do zero
  \item O arquivo de memória (CLAUDE.md etc.) é a memória de longo prazo e carrega sozinho em cada sessão
  \item Bom arquivo = visão + stack + padrões + pastas + convenções; 70 linhas bastam
  \item SDD: escreva a spec antes do código; a spec é a fonte única da verdade e é mais estável
  \item Contexto em quatro camadas: memória do projeto @arquivos histórico instrução imediata
  \item Cinco práticas: atualizar, usar @, chat curto, mostrar o original, limpar
\end{itemize}

\end{tcbsummary}
